\documentclass[12pt]{article}
\usepackage[utf8]{inputenc}
\usepackage[T1]{fontenc}
\usepackage{hyperref}
\title{Munda Language Resources for K-12 Curriculum Development}
\author{Compiled by Claude Sonnet 5 (claude-sonnet-5) for Andr\'as Kornai}
\date{}
\begin{document}
\maketitle

\section{Introduction}
The code unx (``Munda'') is, according to Glottolog, a retired ISO 639-3 entry kept only for bookkeeping. Glottolog notes that Munda [unx] and Mundari [unr] are two entries for what is considered the same language. An archive record at CIIL Mysuru lists the glottocode mund1321 with unx, but I could not confirm mund1321 as a current Glottolog entry. Mundari itself is listed under unr with Glottocode mund1320. This dossier therefore cannot treat unx as a distinct language with its own resources. Any curriculum work should go through the Mundari (unr) entry. I did not verify that any school materials exist.

\section{Primary Educational Resources}
No primer, textbook or school programme specific to unx was found in the search results.

\section{Linguistic and Reference Materials}
\begin{itemize}
\item Glottolog bookkeeping page for Munda [unx]: \url{https://glottolog.org/resource/languoid/iso/unx}
\item Wikipedia, Mundari language (background on the corresponding unr language): \url{https://en.wikipedia.org/wiki/Mundari_language}
\item Wikipedia, Munda languages (family overview): \url{https://en.wikipedia.org/wiki/Munda_languages}
\end{itemize}

\section{Religious and Community Resources}
None found for unx specifically.

\section{Digital and Audio Resources}
\begin{itemize}
\item CIIL Mysuru archive record ``Songs and conversation in Munda'' (it lists the glottocode mund1321 with unx): \url{https://sanchika.ciil.org/items/483b5b65-b804-4cba-9c3e-7e2608163518/full}
\end{itemize}

\section{Curriculum Development Context}
Because unx is retired and overlaps with Mundari (unr), a curriculum effort should be planned under Mundari and with Indian authorities and community bodies. Little else was established by this search, and nothing here should be read as evidence of existing school provision.
\end{document}